\chapter{Factors, Matrices, and Arrays}
\label{chap:factors-matrices-arrays}

\chapterguide
  {You should be comfortable with vectors, lists, indexing, and arithmetic.}
  {create and inspect factors, matrices, and arrays; index their dimensions; perform element-wise arithmetic; and use \texttt{apply()}.}

\section{Factors}

Convert a character vector of regions to a factor:

\begin{lstlisting}[style=dsr]
data <- c(
  "East", "West", "East", "North", "North",
  "East", "West", "West", "West", "East", "North"
)

print(data)
print(is.factor(data))

factor_data <- factor(data)
print(factor_data)
print(is.factor(factor_data))
length(factor_data)
\end{lstlisting}

\texttt{is.factor()} confirms the conversion.

\begin{utpfig}[Inside a factor]{A factor is not a character vector. It stores
integer codes plus a \texttt{levels} attribute holding each label once. The
printed output shows labels, which is why the difference stays hidden until
something behaves unexpectedly.}{fig:factor}
\begin{tikzpicture}[x=1cm,y=1cm]
\node[tagb,anchor=east] at (-0.15,0.55) {what you see};
\foreach \v/\i in {East/0, West/1, East/2, North/3, North/4, East/5}
  \node[cell,minimum width=13mm,font=\tiny\color{UTPNavy}] at (\i*1.42,0.55) {\v};
\node[tag,anchor=west] at (7.6,0.55) {\ldots};

\node[tagb,anchor=east] at (-0.15,-0.55) {what is stored};
\foreach \v/\i in {1/0, 3/1, 1/2, 2/3, 2/4, 1/5}
  \node[cellalt,minimum width=13mm] at (\i*1.42,-0.55) {\v};
\node[tag,anchor=west] at (7.6,-0.55) {\ldots};

\node[tagb,anchor=east] at (-0.15,-1.7) {\texttt{levels}};
\foreach \v/\i in {{1: East}/0, {2: North}/1, {3: West}/2}
  \node[cellout,minimum width=17mm,font=\tiny\color{UTPNavy}] at (\i*1.85,-1.7) {\texttt{\v}};
\node[tag,anchor=west,text width=4.4cm] at (5.4,-1.7)
  {alphabetical by default --- not the order they appeared};
\draw[rarrowmuted] (0,-1.32) -- (0,-0.92);
\end{tikzpicture}
\end{utpfig}

\begin{pitfall}
Two consequences follow from the codes. \texttt{as.numeric()} on a factor
returns the codes, not the labels: on a factor of years,
\texttt{as.numeric(f)} gives \texttt{1 2 3}, not \texttt{2021 2022 2023}. Go
through \texttt{as.character()} first. And a value that is not already a level
cannot be assigned --- R inserts \texttt{NA} and warns.
\end{pitfall}

\subsection{Accessing and modifying factor values}

\begin{lstlisting}[style=dsr]
data <- factor(c(
  "East", "West", "East", "North", "North",
  "East", "West", "West", "West", "East", "North"
))

data[3]
data[3] <- "NorthWest"
\end{lstlisting}

Assigning a value outside the existing levels produces \texttt{NA} with a warning. Add the level first if it is valid.

\subsection{Changing level order}

\begin{lstlisting}[style=dsr]
data <- c(
  "East", "West", "East", "North", "North",
  "East", "West", "West", "West", "East", "North"
)

factor_data <- factor(data)
new_order_data <- factor(
  factor_data,
  levels = c("East", "West", "North")
)
print(new_order_data)
\end{lstlisting}

The \texttt{levels} argument controls category order independently of observation order.

\subsection{Generating factor levels}

\begin{lstlisting}[style=dsr]
v <- gl(3, 4, labels = c("Tampa", "Seattle", "Boston"))
print(v)
\end{lstlisting}

\texttt{gl()} generates repeated factor levels.

\section{Matrices}

A matrix is one atomic vector with a \texttt{dim} attribute attached. That
single fact explains both the fill order and why a matrix cannot mix types.

\begin{utpfig}[Matrix fill order]{\texttt{matrix(1:6, nrow = 2)} against the same
call with \texttt{byrow = TRUE}. R fills down the first column before starting
the second unless told otherwise, which is the reverse of how the numbers read
on the page.}{fig:matrix-fill}
\begin{tikzpicture}[x=1cm,y=1cm]
\node[tagb,anchor=east] at (-0.2,0.9) {\texttt{1:6}};
\foreach \v/\i in {1/0, 2/1, 3/2, 4/3, 5/4, 6/5}
  \node[cell] at (\i*0.9,0.9) {\v};

% column-major
\node[tagb] at (1.3,-0.15) {default (column-major)};
\foreach \v/\c/\r in {1/0/0, 3/1/0, 5/2/0, 2/0/1, 4/1/1, 6/2/1}
  \node[cellalt] at (0.45+\c*0.9,-0.75-\r*0.75) {\v};
\draw[rarrowmuted] (0.45,-0.55) -- (0.45,-1.7);
\node[tag,anchor=west] at (0.62,-1.75) {fills down first};

% row-major
\node[tagb] at (5.6,-0.15) {\texttt{byrow = TRUE}};
\foreach \v/\c/\r in {1/0/0, 2/1/0, 3/2/0, 4/0/1, 5/1/1, 6/2/1}
  \node[cellout] at (4.75+\c*0.9,-0.75-\r*0.75) {\v};
\draw[rarrowmuted] (4.4,-0.75) -- (6.15,-0.75);
\node[tag,anchor=west] at (4.4,-1.75) {fills across first};
\end{tikzpicture}
\end{utpfig}

A matrix stores one data type in rows and columns. \texttt{byrow} controls the fill direction:

\begin{lstlisting}[style=dsr]
M <- matrix(c(3:14), nrow = 4, byrow = TRUE)
print(M)

N <- matrix(c(3:14), nrow = 4, byrow = FALSE)
print(N)
\end{lstlisting}

\subsection{Naming rows and columns}

\begin{lstlisting}[style=dsr]
rownames <- c("row1", "row2", "row3", "row4")
colnames <- c("col1", "col2", "col3")

P <- matrix(
  c(3:14),
  nrow = 4,
  byrow = TRUE,
  dimnames = list(rownames, colnames)
)
print(P)
\end{lstlisting}

Dimension names improve readability and indexing.

\subsection{Matrix indexing}

\begin{lstlisting}[style=dsr]
print(P[1, 3])  # row 1, column 3
print(P[4, 2])  # row 4, column 2
print(P[2, ])   # all columns from row 2
print(P[, 3])   # all rows from column 3
\end{lstlisting}

Matrix indices follow \texttt{[row, column]}.

\section{Element-wise matrix arithmetic}

Ordinary operators act element by element on conformable matrices:

\begin{lstlisting}[style=dsr]
matrix1 <- matrix(c(3, 9, -1, 4, 2, 6), nrow = 2)
matrix2 <- matrix(c(5, 2, 0, 9, 3, 4), nrow = 2)

print(matrix1 + matrix2)
print(matrix1 - matrix2)
print(matrix1 * matrix2)
print(matrix1 / matrix2)
\end{lstlisting}

\begin{keyidea}
\texttt{*} performs element-wise multiplication; true matrix multiplication uses \texttt{\%*\%}.
\end{keyidea}

\section{Arrays}

Arrays extend matrices to additional dimensions.

\subsection{Constructing an array}

\begin{lstlisting}[style=dsr]
vector1 <- c(5, 9, 3)
vector2 <- c(10, 11, 12, 13, 14, 15)

result <- array(
  c(vector1, vector2),
  dim = c(3, 3, 2)
)
print(result)
\end{lstlisting}

The dimensions specify 3 rows, 3 columns, and 2 layers. Because only 9 values are supplied for 18 positions, R recycles them once.

\subsection{Naming array dimensions}

\begin{lstlisting}[style=dsr]
column.names <- c("COL1", "COL2", "COL3")
row.names <- c("ROW1", "ROW2", "ROW3")
matrix.names <- c("Matrix1", "Matrix2")

result <- array(
  c(vector1, vector2),
  dim = c(3, 3, 2),
  dimnames = list(row.names, column.names, matrix.names)
)
print(result)
\end{lstlisting}

\subsection{Accessing array elements and slices}

\begin{lstlisting}[style=dsr]
print(result[3, , 2])  # third row of second matrix
print(result[1, 3, 1]) # row 1, column 3, first matrix
print(result[, , 2])   # entire second matrix
\end{lstlisting}

Array indices follow \texttt{[row, column, layer]}.

\section{Array manipulation and \texttt{apply()}}

\texttt{apply()} runs a function across selected dimensions:

\begin{lstlisting}[style=dsr]
new.array <- array(
  c(vector1, vector2),
  dim = c(3, 3, 2)
)

result <- apply(new.array, c(1), sum)
print(result)
\end{lstlisting}

Margin \texttt{1} calculates one sum per row across columns and layers.

\begin{pitfall}
Lab 4a defines \texttt{vector3} and \texttt{vector4} but constructs \texttt{array2} from the earlier vectors. Confirm the intended inputs before reproducing that step.
\end{pitfall}

\begin{quickcheck}
\begin{enumerate}
  \item Which function converts the region vector into a factor?
  \item What does \texttt{byrow = TRUE} change when constructing a matrix?
  \item In \texttt{P[4,2]}, which coordinate is the row and which is the column?
  \item How many dimensions are specified by \texttt{dim = c(3,3,2)}?
  \item What does \texttt{result[,,2]} select from a three-dimensional array?
\end{enumerate}
\end{quickcheck}

\begin{chapterrecap}
\begin{itemize}
  \item Factors store categorical values through defined levels.
  \item Matrices and arrays use dimension-based indexing and may carry names.
  \item Ordinary arithmetic is element-wise; \texttt{apply()} summarizes selected dimensions.
\end{itemize}
\end{chapterrecap}
